\section{A linear bound on the low-energy state count}
\label{sec:counting}

Let $\mathcal N_{N,w}(u)$ count the transfer energies at circumference $w$
that are at most $u$, with multiplicity and including the vacuum.
The visibility estimate from Section~\ref{sec:visibility} now allows us to
turn the trace bounds into a count below the two-excitation threshold.

\begin{proposition}\label{cou:linear}
There is a constant $C$ such that, for every sufficiently large dyadic
circumference $w=w_k$,
\begin{equation}\label{cou:linear-eq}
 \limsup_{N\to\infty}\mathcal N_{N,w}(u_0)\le Cw,
 \qquad u_0=\frac32m_*.
\end{equation}
\end{proposition}

Set $u_1=13m_*/8$. The intermediate goal is the estimate
\[
 \mathcal N_{N,2w}(y)
 \le 2(1+\varepsilon_w)\mathcal N_{N,w}(y+b_w),
 \qquad u_0\le y\le u_1,
\]
with positive errors $b_w,\varepsilon_w$ summable over dyadic widths.
For each sufficiently large fixed $w$, this estimate will hold for all
sufficiently large $N$, with the cutoff threshold allowed to depend on
$w$. Summability of the energy shifts keeps the iterated threshold inside
the stated interval once the base width is large enough; the factors
$1+\varepsilon_w$ have a finite product. Thus the factor $2$ at each
doubling yields the desired linear growth in circumference.

To prove the estimate, we compare the heat trace at circumference $2w$
with the trace of two independent copies at circumference $w$. The
energies in the latter system are sums $E+E'$. If such a sum lies below
twice the gap, at least one component is the vacuum energy. A nonnegative
polynomial approximation to an energy cutoff makes this observation
quantitative without assuming convergence of individual eigenvalues.

\subsection{Exchanging time and circumference}

Rotation of the partition function gives the exact identity
\begin{equation}\label{cou:exchange}
 e^{t g_N(w)}R_{N,w}(t)
   =e^{w g_N(t)}R_{N,t}(w),
\end{equation}
where $g_N$ is the pressure correction from Section~\ref{sec:transfer}.
For a large circumference $w$, set
\[
 D=\lfloor w^{1/2}\rfloor,\qquad t_0=\lfloor D^{1/2}\rfloor.
\]
Every integer physical time is an allowed lattice length. Uniformly in the
cutoff and in the integers $t_0\le t\le D$, we claim that
\begin{equation}\label{cou:heat-product}
 \bigl|R_{N,2w}(t)-R_{N,w}(t)^2\bigr|\le Ce^{-cw}.
\end{equation}
Indeed, applying \eqref{cou:exchange} at $w$ and $2w$ yields
\begin{align*}
 \log R_{N,2w}(t)-2\log R_{N,w}(t)
  ={}&t\bigl(2g_N(w)-g_N(2w)\bigr)\\
    &+\log R_{N,t}(2w)-2\log R_{N,t}(w).
\end{align*}
The first term is exponentially small in $w$ by \eqref{tra:trace}.
The uniform positive gap at all sufficiently large circumferences,
Corollary~\ref{tra:gap}, gives a constant $c_*>0$ such that
\[
 s_{N,t}(2w)\le s_{N,t}(w)
 \le e^{-c_*(w-t)}s_{N,t}(t)\le Ce^{-c'w}.
\]
Here $s_{N,t}(t)$ is uniformly bounded by \eqref{tra:trace}.
Thus the difference of logarithms is exponentially small.
The same trace estimate bounds both $R_{N,2w}(t)$ and $R_{N,w}(t)^2$
uniformly on this time range: their logarithms are at most a polynomial
in $w$ times $e^{-ct_0}$. Exponentiating proves
\eqref{cou:heat-product}.

\subsection{A positive approximation to an energy cutoff}

For $u_0\le y\le u_1$, define
\[
 b_w=D^{-1/8},\qquad
 j_0=\left\lceil D e^{-(y+b_w/2)}\right\rceil,
 \qquad
 Q_{w,y}(x)=\sum_{j=j_0}^{D}\binom Dj x^j(1-x)^{D-j}.
\]
The polynomial is the upper tail of a binomial law, so
$0\le Q_{w,y}(x)\le1$ on $[0,1]$. Its transition occurs between energies
$y$ and $y+b_w$ when $x=e^{-E}$.

\begin{lemma}\label{cou:polynomial}
For all sufficiently large $w$, uniformly in $y\in[u_0,u_1]$,
\begin{align}
 Q_{w,y}(e^{-E})&\ge1-Ce^{-cD^{3/4}}
       &&(0\le E\le y),\label{cou:poly-low}\\
 Q_{w,y}(e^{-E})&\le Ce^{-cD^{3/4}}e^{-t_0E}
       &&(E\ge y+b_w).\label{cou:poly-high}
\end{align}
Every nonzero monomial of $Q_{w,y}$ has degree between $t_0$ and $D$,
and the sum of the absolute values of its coefficients is at most $3^D$.
\end{lemma}

\begin{proof}
The threshold $j_0/D$ differs from $e^{-y}$ and $e^{-(y+b_w)}$ by at
least $c b_w$ on their respective sides. This follows from the mean value
theorem, uniformly on the fixed energy interval, since the rounding error
is at most $1/D=o(b_w)$. Hoeffding's binomial tail bound
\cite[Theorem~1]{Hoeffding1963} therefore gives \eqref{cou:poly-low} and
\eqref{cou:poly-high} without the factor $e^{-t_0E}$, because
$Db_w^2=D^{3/4}$.

To retain a summable bound over all energies, that additional factor is
essential. Choose $c_1>0$ such that $j_0\ge c_1D$ for every allowed $y$
and all large $w$. Fix $E_*>4\log(2)/c_1+1$.
For $E\le E_*$, the loss $e^{t_0E_*}$ is absorbed by decreasing the
constant in the exponent $-cD^{3/4}$, since $t_0=O(D^{1/2})$.
For $E\ge E_*$, the binomial formula gives
\[
 Q_{w,y}(e^{-E})\le 2^D e^{-j_0E}.
\]
For large $w$, $t_0\le c_1D/2$, and consequently
\[
 2^D e^{-(j_0-t_0)E}
 \le \exp\!\left(-\frac{c_1}{4}DE\right)
 \le Ce^{-cD^{3/4}}.
\]
This proves \eqref{cou:poly-high} for every $E\ge y+b_w$.

The smallest possible degree in the expanded polynomial is $j_0$, which
is at least $t_0$ for large $w$. Finally its coefficient sum in absolute
value is bounded by
\[
 \sum_{j=j_0}^D\binom Dj 2^{D-j}\le3^D.
\]
\end{proof}

Applying \eqref{cou:heat-product} to each monomial gives
\begin{equation}\label{cou:poly-trace}
 \left|
   \sum_{E\text{ at }2w}Q_{w,y}(e^{-E})
   -\sum_{E,E'\text{ at }w}Q_{w,y}(e^{-(E+E')})
 \right|\le Ce^{-c'w}.
\end{equation}
All sums count multiplicities. They converge absolutely, since the
smallest degree is positive and the relevant heat traces are finite.
The coefficient bound costs at most $3^D$, which is absorbed by
$e^{-cw}$ because $D=O(w^{1/2})$.

\paragraph{From polynomial traces to counts.}
By the visibility estimate \eqref{vis:gap}, for each sufficiently large
fixed dyadic $w$ and all sufficiently large $N$ depending on $w$,
\[
 2e_N(w)\ge\frac74m_*.
\]
Increase the lower bound on $w$ so that $u_1+b_w<7m_*/4$. Any pair in
\eqref{cou:poly-trace} with $E+E'<y+b_w$ then has a zero component.
There are at most $2\mathcal N_{N,w}(y+b_w)$ such pairs. For the
remaining pairs, Lemma~\ref{cou:polynomial} bounds their total contribution
by
\[
 Ce^{-cD^{3/4}}\sum_{E,E'}e^{-t_0(E+E')}
 =Ce^{-cD^{3/4}}R_{N,w}(t_0)^2
 \le Ce^{-cD^{3/4}}.
\]
Using \eqref{cou:poly-low} for energies at $2w$ not exceeding $y$, and
absorbing the additive errors into $\mathcal N_{N,w}\ge1$, proves
\begin{equation}\label{cou:recursion}
 \mathcal N_{N,2w}(y)
 \le2\bigl(1+Ce^{-cD^{3/4}}\bigr)
       \mathcal N_{N,w}(y+b_w),
 \qquad u_0\le y\le u_1.
\end{equation}
The constants are independent of $N,w,y$; for each fixed $w$, the
inequality holds for all sufficiently large $N$.

\subsection{Dyadic iteration}

\begin{proof}[Proof of Proposition~\ref{cou:linear}]
Choose a fixed dyadic base width $w_{k_0}$ large enough for the preceding
estimates and so that
\[
 \sum_{k=k_0}^{\infty}b_{w_k}<u_1-u_0.
\]
This is possible because $b_{w_k}=O(w_k^{-1/16})$.
The product of the factors $1+Ce^{-c\lfloor\sqrt{w_k}\rfloor^{3/4}}$
over the same range is finite. At the base width,
\[
 \mathcal N_{N,w_{k_0}}(u_1)
 \le e^{w_{k_0}u_1}R_{N,w_{k_0}}(w_{k_0})\le C_0,
\]
uniformly in the cutoff.

Fix a target $w_k$ with $k>k_0$. Apply \eqref{cou:recursion} backwards
along the finite chain $w_k,w_{k-1},\ldots,w_{k_0}$, starting at threshold
$u_0$. The accumulated threshold shifts remain below $u_1-u_0$, so every
application is valid. For all sufficiently large $N$ on this finite
chain, we obtain
\[
 \mathcal N_{N,w_k}(u_0)
 \le C_0\,2^{k-k_0}
       \prod_{h=k_0}^{k-1}
         \left(1+Ce^{-c\lfloor\sqrt{w_h}\rfloor^{3/4}}\right)
 \le Cw_k.
\]
Taking $\limsup_{N\to\infty}$ proves the assertion. The cutoff threshold
may depend on the target width; no simultaneous limit in $N$ and $w$
has been used.
\end{proof}
